\section{Introduction}
\label{sec:introduction}

Language models can generate RTL that passes executable benchmarks
~\cite{liu2023verilogeval,lu2024rtllm,pinckney2025revisiting}, but functional
acceptance alone says little about the quality of the implemented circuit.
Modules that compute the same outputs can differ substantially in their critical
paths, resource mappings, and achievable clock frequencies. A model trained
only to pass a testbench has no direct incentive to prefer the faster
implementation. Sampling several alternatives and selecting the fastest can
address this problem, although it requires repeated electronic design automation
(EDA) runs at inference time.

Recent work uses synthesis feedback, reinforcement learning, evolutionary
search, and per-design adaptation to improve RTL
~\cite{ppartl2025,chipseek2025,yao2024rtlrewriter,poet2026,
evolve2026,zhou2026alphartl,drrtl2026}. These approaches make it important to
distinguish the quality of a selected circuit from the expected quality of a
single policy draw. Synthesis-based evaluation also finds a gap between
single-attempt and best-of-several results~\cite{fu2026synthesis}. A learned
physical reward offers a way to guide training without implementing every
candidate, but optimization can shift the generated code toward regions where
the predictor is less reliable~\cite{gao2023scaling}. Its training scores must
therefore be assessed against independent physical measurements.

We study whether training with a physical proxy improves expected physical
utility at a fixed draw budget in a controlled streaming-accelerator domain.
Both functional and implementation failures contribute zero, so the endpoint
reflects the full distribution of generated outputs. Matched reward controls
test whether physical feedback adds value beyond correctness pressure alone.
As shown in Fig.~\ref{fig:training-boundary}, domain SFT initializes the
generator, an executable oracle gates the physical reward, and group-relative
updates change candidate probabilities. Frozen checkpoints are then evaluated
through independent post-route implementation. The contribution lies in this
distribution-level evaluation and controlled reward comparison; the policy
update and EDA feedback build on established methods.

\begin{figure*}[!t]
\centering
\includegraphics[width=\textwidth]{fig01_training_evidence_boundary_v3}
\caption{Implemented training and physical-evaluation workflow.
(a) Domain SFT and the physical predictor use training-only resources.
(b) Fresh policy samples receive group-relative updates with a frozen SFT
reference; flat groups do not advance the update counter.
(c) The RF reward uses canonicalized structural features; MLP and
correctness-only rewards are separate control runs. Red branches assign zero
to rejected candidates.
(d) Frozen checkpoints are evaluated using the executable oracle and post-route
implementation. Original multiplicity and the full draw denominator are
restored; incorrect and implementation-failed draws are not discarded.
Dashed paths denote feedback, references, or separately run controls.}
\label{fig:training-boundary}
\end{figure*}

The primary test is a separately frozen
\studyclaim{StudyTwoDesignCount}-design study spanning interpolation and
extrapolation across the retained accelerator generators.  Two policy-training
seeds sharing one SFT checkpoint and a canonicalization-aware random-forest (RF) reward raise
penalized equal-sample \fmax{} from
\studyclaim{StudyTwoOverallSftPenalized} to
\studyclaim{StudyTwoOverallRfPenalized}~MHz.  Their seed-averaged paired effect
is \studyclaim{StudyTwoOverallRfSftGain}~MHz with interval
[\studyclaim{StudyTwoOverallRfSftCiLow},
\studyclaim{StudyTwoOverallRfSftCiHigh}]~MHz; each training seed is positive,
and the interpolation and extrapolation intervals both exclude zero.  The
optimized policy's correctness remains close to the supervised fine-tuned
(SFT) baseline, and the failure penalty already charges the endpoint for the
difference.

The matched controls help explain the gain.  The repaired RF reward
exceeds the original multilayer-perceptron (MLP) reward by
\studyclaim{StudyTwoOverallRfMlpGain}~MHz under the same tasks and sampling
budget.  Correctness-only training instead raises validity to
\studyclaim{StudyTwoCorrectnessOnlyCorrectness}\% but lowers penalized \fmax{}
relative to SFT in its single preregistered seed.  The observed physical gain
therefore does not arise from higher functional yield.  A separate completed
setup-closure extension repeats the RF and correctness-only physical directions
in additional policy-training seeds, while retaining the shared-SFT limitation.
A separate Qwen-Coder~\cite{hui2024qwen25coder} experiment improves utility on all
\claim{QwenDesignCount} tested designs. Conditional resource measurements,
paired board checks, and general-RTL evaluation describe the engineering costs
and limits of these gains.

Together, these experiments establish a controlled physical improvement on
the frozen split and identify its engineering limits. Family and sensitivity
analyses show where the gains are weaker, while the original-reward and
correctness-only arms distinguish physical improvement from increased
functional yield. The comparison supports the RF reward package as a whole;
it does not isolate predictor architecture from canonical features. Throughout
the evaluation, the fixed-draw endpoint preserves sample multiplicity and
assigns zero to every incorrect or implementation-failed draw. Post-route
measurements determine the reported outcome, with source-level RTL analysis,
conditional PPA, setup checks, board measurements, and general-RTL evaluation
providing context for its practical value.

The candidate analysis suggests that much of the gain comes from making fast
implementations more likely. Supervision already supplies several passing
forms, and optimization concentrates probability on a smaller set of outputs.
The sealed support diagnostic nevertheless contains exceptions, so the finite
SFT sample cannot account for every observed gain. These results establish a
change in the sampled distribution without resolving whether optimization
produces new circuit structures. They also explain why the benefit depends
on the implementations available to the generator and on the reward's ability
to distinguish their physical quality.
